\chapter{Nose cancer}
\index{Nose cancer}

\section*{Definition and Overview}
Nose cancer refers to malignant neoplasms arising in the nasal cavity (the air-filled space behind the nose) and the paranasal sinuses (the air-filled cavities within the bones surrounding the nasal passages). These cancers are rare and typically originate in the squamous cells that line these passages, though they can also arise from salivary gland cells (such as adenoid cystic carcinoma), mucosal melanocytes, or neuroectodermal tissues (such as esthesioneuroblastoma). These malignancies present significant diagnostic challenges due to their anatomical complexity and the tendency for early symptoms to mimic common benign conditions, such as chronic sinusitis or allergic rhinitis, which often results in diagnosis at an advanced stage.

\section*{Epidemiology}
Nasal cavity and paranasal sinus cancers are rare, representing less than 1\% of all malignancies worldwide and approximately 3\% to 5\% of all head and neck cancers. They occur more commonly in males than females, with a ratio of approximately 2:1, and are most frequently diagnosed in individuals aged between 50 and 70 years. In many countries, the incidence remains low but stable, with a strong association with occupational exposures. Wood dust, leather dust, nickel, and chromium compounds are well-documented risk factors, particularly for adenocarcinoma of the ethmoid sinuses. Tobacco smoke and, to a lesser extent, human papillomavirus (HPV) infection are also implicated in squamous cell carcinomas of the nasal cavity.

\section*{Aetiology and Pathophysiology}
The development of nose cancer involves a multi-step process of genetic alterations in the mucosal lining of the nasal cavity and paranasal sinuses, driven largely by chronic exposure to inhaled carcinogens.
\begin{itemize}
    \item \textbf{Occupational carcinogens:} Chronic inhalation of wood dust, leather dust, nickel dust, and formaldehyde causes persistent mucosal irritation, squamous metaplasia, and eventual malignant transformation, especially in the ethmoid sinuses.
    \item \textbf{Tobacco smoking:} Active smoking exposes the nasal epithelium to inhaled carcinogens, increasing the risk of squamous cell carcinoma.
    \item \textbf{Human papillomavirus (HPV):} Certain high-risk strains of HPV, particularly HPV-16, are associated with a subset of squamous cell carcinomas and inverted papillomas that undergo malignant progression.
    \item \textbf{Chronic inflammation:} Long-standing inflammatory conditions, though less directly causative, may create a microenvironment conducive to cellular damage and mutation.
    \item \textbf{Genetic mutations:} Accumulation of mutations in tumour suppressor genes (such as TP53) and activation of oncogenes disrupt normal cell cycle regulation, prompting uncontrolled cellular proliferation.
\end{itemize}

\section*{Clinical Features}
Symptoms of nose cancer are often unilateral and progressive, distinguishing them from bilateral benign conditions.
\begin{itemize}
    \item \textbf{Nasal obstruction:} Persistent, progressive blockage of one side of the nose that does not respond to nasal sprays or decongestants.
    \item \textbf{Epistaxis:} Frequent, recurrent nosebleeds, typically unilateral, which may start as blood-tinged nasal discharge.
    \item \textbf{Nasal discharge:} Persistent anterior rhinorrhoea or post-nasal drip, which is often purulent, foul-smelling, or blood-stained.
    \item \textbf{Facial pain:} Pain, headache, or pressure in the cheek, forehead, or around the eyes.
    \item \textbf{Sensory disturbances:} Decreased sense of smell (hyposmia) or complete loss of smell (anosmia).
    \item \textbf{Ocular symptoms:} Double vision (diplopia), bulging of the eye (proptosis), tearing (epiphora), or visual loss, occurring as the tumour invades the orbit.
    \item \textbf{Oral and dental symptoms:} Loosening of upper teeth, pain in the jaw, or difficulty opening the mouth (trismus) due to infiltration of the maxillary sinus floor.
    \item \textbf{Lymphadenopathy:} Swelling or painless lumps in the neck, indicating regional lymph node metastasis.
\end{itemize}

\section*{Differential Diagnosis}
Because symptoms are non-specific, it is critical to distinguish nose cancer from benign inflammatory and structural disorders.
\begin{itemize}
    \item \textbf{Chronic rhinosinusitis:} Bilateral inflammatory condition of the nasal passages and sinuses, presenting with congestion, facial pressure, and discharge, but typically responding to medical therapy.
    \item \textbf{Nasal polyps:} Benign mucosal growths that are usually bilateral and associated with asthma or allergies, unlike unilateral malignant lesions.
    \item \textbf{Inverted papilloma:} A benign but locally aggressive epithelial tumour of the nasal cavity that carries a significant risk of malignant transformation.
    \item \textbf{Granulomatosis with polyangiitis:} A systemic autoimmune vasculitis that causes necrotising granulomas, nasal ulceration, and septal perforation.
    \item \textbf{Allergic rhinitis:} A common allergic response causing bilateral congestion, sneezing, and clear rhinorrhoea, lacking red-flag signs like unilateral bleeding.
    \item \textbf{Foreign body obstruction:} Unilateral nasal discharge and blockage, primarily seen in paediatric patients.
\end{itemize}

\section*{When to Seek Medical Care}
Individuals experiencing persistent, unilateral nasal symptoms—such as obstruction, recurrent bleeding, or facial pain—that last longer than three to four weeks should seek medical evaluation from a general practitioner or an ear, nose, and throat (ENT) specialist. Prompt assessment is essential to rule out malignancy and prevent advanced disease progression.

\textbf{Contact your local emergency services for:}
\begin{itemize}
    \item Severe, uncontrollable nosebleeds (epistaxis) causing respiratory distress or hemodynamic instability
    \item Sudden loss of vision, double vision, or inability to move the eye
    \item Severe, sudden-onset headache or facial swelling accompanied by high fever or confusion
    \item Difficulty breathing, airway obstruction, or stridor
\end{itemize}

\section*{Diagnosis and Investigations}
A comprehensive diagnostic workup is required to visualise the lesion, obtain a tissue diagnosis, and determine the extent of local and distant spread.
\begin{itemize}
    \item \textbf{Nasal endoscopy:} A flexible or rigid endoscope is inserted into the nasal cavity to visualise the mucosa, identify abnormal growths, and locate the site of bleeding.
    \item \textbf{Biopsy:} Retrieval of a tissue sample, often during endoscopy under local or general anaesthetic, for histopathological examination to confirm malignancy and type.
    \item \textbf{CT scan:} High-resolution computed tomography of the paranasal sinuses and skull base to evaluate bone erosion and structural involvement.
    \item \textbf{MRI scan:} Magnetic resonance imaging to assess soft tissue invasion, perineural spread, and extension into the orbit or brain.
    \item \textbf{PET-CT scan:} Positron emission tomography combined with CT, utilised to detect regional lymph node metastases and distant haematogenous spread.
    \item \textbf{Fine-needle aspiration (FNA) biopsy:} Performed on palpable neck nodes to check for regional metastasis.
\end{itemize}

\section*{Management}
Treatment is highly individualised and determined by a multidisciplinary head and neck oncology team, depending on the stage, location, histology, and patient's functional status.
\begin{itemize}
    \item \textbf{Surgery:} Surgical resection remains the cornerstone of treatment, ranging from minimally invasive endoscopic resection to open craniofacial resection for advanced tumours.
    \item \textbf{Radiotherapy:} External beam radiation therapy (such as intensity-modulated radiation therapy) used as primary treatment for inoperable tumours, or postoperatively to eradicate microscopic residual disease.
    \item \textbf{Chemotherapy:} Systemic chemotherapy (e.g., cisplatin) administered concurrently with radiation (chemoradiotherapy) for advanced, high-grade, or unresectable disease.
    \item \textbf{Targeted therapy and immunotherapy:} Used in recurrent or metastatic cases, targeting specific molecular pathways or enhancing the immune response against cancer cells.
    \item \textbf{Rehabilitation and supportive care:} Speech pathology, nutritional support, and prosthetic reconstruction (obturators) to restore swallowing, speech, and facial appearance after extensive surgery.
\end{itemize}

\section*{Complications}
\begin{itemize}
    \item Local invasion into the brain, skull base, or orbit, leading to visual loss, meningitis, or cerebrospinal fluid (CSF) leaks
    \item Severe cosmetic and structural deformity of the midface resulting from disease progression or surgical resection
    \item Osteoradionecrosis (bone death due to radiation) of the maxilla or skull base
    \item Severe xerostomia (dry mouth), mucositis, and swallowing dysfunction secondary to radiotherapy
    \item Sensory loss, including permanent anosmia (loss of smell) and gustatory impairment
    \item Systemic complications of chemotherapy, including nephrotoxicity, ototoxicity, and neutropenic sepsis
\end{itemize}

\section*{Prognosis and Outcomes}
The prognosis for nose cancer is heavily dependent on the stage at presentation, the specific histological subtype, and the adequacy of surgical margins. Early-stage, localised tumours (Stage I and II) carry a favourable five-year survival rate of approximately 60\% to 80\% following successful surgical resection and adjuvant therapy. However, because many patients present with advanced disease (Stage III and IV) due to the insidious onset of symptoms, the overall five-year survival rate falls to approximately 35\% to 50\%. Tumours that invade the skull base, orbit, or brain, or those demonstrating regional or distant metastases, carry a significantly poorer prognosis and require aggressive multimodal management and long-term surveillance.

\section*{Prevention and Self-Care}
\begin{itemize}
    \item Minimise occupational exposure to wood dust, leather dust, formaldehyde, and heavy metals by using appropriate personal protective equipment (PPE) and proper ventilation.
    \item Avoid tobacco smoking and seek support to quit, as smoking increases the risk of head and neck cancers.
    \item Maintain good hygiene in industrial workplaces and undergo regular health screenings if exposed to known nasopharyngeal carcinogens.
    \item Adhere strictly to post-treatment surveillance schedules, including regular endoscopies and imaging, to detect early recurrence.
    \item Manage dry nasal passages using prescribed saline rinses and humidifiers, especially during and after radiotherapy.
\end{itemize}

\section*{Sources and Further Reading}
The following reputable organisations provide further information and clinical guidelines on nose cancer:
\begin{itemize}
    \item Cancer Council Australia. \textit{Information on head, neck, and nasal cancers.} https://www.cancer.org.au/
    \item Mayo Clinic. \textit{Comprehensive overview of nasal and paranasal sinus cancers.} https://www.mayoclinic.org/
    \item NHS. \textit{Patient guide to nasal and sinus cancer symptoms and treatment.} https://www.nhs.uk/
    \item MedlinePlus. \textit{Nasal cancer information, resources, and clinical trials.} https://medlineplus.gov/
    \item National Cancer Institute. \textit{Detailed professional and patient resources on paranasal sinus and nasal cavity cancer.} https://www.cancer.gov/
\end{itemize}
